\section{Resident State}
Resident state is the memory floor that exists before we talk about
phase-specific activations or temporary workspace. A useful beginner rule is:

\begin{quote}
If the rank still owns the object even while no kernel is actively running,
that object belongs in the resident floor.
\end{quote}

\subsection{The Basic Byte Rule}
Every resident tensor follows the same local accounting rule:
\[
\text{local bytes} = \text{local elements} \times \text{bytes per element}.
\]

The only real question is what ``local'' means. Tensor parallelism reduces the
visible width of some tensors. ZeRO may further shard trainable tensors across
the data-parallel group. Once those local element counts are known, the byte
computation itself is straightforward.

The two local-count rules that show up throughout the estimator are:
\begin{align*}
n_{\mathrm{local,tp}} &\approx \dfrac{n}{d_{\mathrm{tp}}}, \\
n_{\mathrm{local,zero}} &\approx \dfrac{n_{\mathrm{local,tp}}}{D_{\mathrm{P}}}.
\end{align*}
In words: first shard by tensor parallel if that tensor is TP-visible, then
shard again by the data-parallel group if the ZeRO stage owns that object in
shards.
The implementation still rounds to whole elements, but the rounded form does
not add much intuition here, so this report shows the simpler proportional
version.

\subsection{Parameters}
Parameter bytes are the cleanest part of the estimator.

\paragraph{Full fine-tuning.}
All base-model parameters are trainable and therefore matter for:
\begin{itemize}[leftmargin=1.5em]
\item resident parameter memory,
\item gradient memory,
\item optimizer state,
\item and, if enabled, master weights.
\end{itemize}

\paragraph{LoRA.}
The frozen base model is still resident because the forward pass needs it, but
only the synthetic adapter parameters are trainable. That distinction matters:
LoRA still pays for the base weights in the resident floor, but it does not pay
for full-model gradients or full-model optimizer state.

\paragraph{ZeRO stages.}
The most important resident distinction is:
\begin{itemize}[leftmargin=1.5em]
\item \texttt{zero2}: parameters remain replicated, even though gradients and
optimizer state are sharded.
\item \texttt{zero3}: trainable parameters are also sharded across the
data-parallel group.
\end{itemize}

At a high level, the resident parameter term is:
\[
M_{\mathrm{params}} =
M_{\mathrm{base\ weights}}
+ M_{\mathrm{adapter\ weights}}.
\]
For full fine-tuning the adapter term is zero. For LoRA the base weights still
exist, but the adapter term is the one that later drives trainable-state
objects such as gradients and optimizer state.

\subsection{Gradients}
Gradient bytes are built from the trainable set, not from the entire parameter
set. This is why LoRA can dramatically reduce gradient memory even though the
base model remains resident.

For distributed training, the estimator treats gradients as:
\begin{itemize}[leftmargin=1.5em]
\item replicated for single-GPU and DDP,
\item sharded for ZeRO-2 and ZeRO-3.
\end{itemize}

That choice is structural. It reflects where gradients live at steady state,
not where they may briefly appear during a communication window.

Written compactly:
\[
M_{\mathrm{grads}} =
\text{local trainable elements} \times \text{gradient bytes per element}.
\]

\subsection{Optimizer State}
Optimizer state is also built from the trainable set. Different optimizers own
different amounts of state:
\begin{itemize}[leftmargin=1.5em]
\item Adam and AdamW keep two auxiliary tensors per trainable tensor.
\item SGD with no momentum keeps essentially none.
\item RMSProp and Adagrad keep one state tensor.
\item Adafactor is special because its state can be factored instead of dense.
\end{itemize}

This is one of the places where the estimator is more principled than a simple
``optimizer multiplier'': it does not assume all optimizers behave the same.
The optimizer family changes the actual state objects.

The compact optimizer-state equation is:
\[
M_{\mathrm{opt}} =
\sum_{p \in \mathcal{P}_{\mathrm{train}}}
\rho_{\mathrm{opt}}(p)\, n_{\mathrm{local}}(p)\, b_o^\star,
\]
where $\rho_{\mathrm{opt}}(p)$ is the optimizer-state ratio for parameter $p$.

\subsection{Master Weights}
Master weights are optional fp32 copies used by some training setups. When they
exist, they are part of the resident floor because they survive across the
entire training step.

\subsection{Runtime Support}
Not every byte on a GPU belongs to a model tensor. CUDA, the caching allocator,
NCCL, and DeepSpeed all keep their own support state. The estimator therefore
models runtime support separately from tensor memory.

The current model exposes four named support components:
\begin{itemize}[leftmargin=1.5em]
\item CUDA context support,
\item allocator-pool support,
\item NCCL support,
\item DeepSpeed support.
\end{itemize}

If the caller wants full control, the estimator also allows one explicit
runtime-support override. The important design choice is not the exact default
numbers. It is that these bytes are now represented as named runtime objects
instead of being hidden inside one opaque per-mode floor.

That term is:
\[
M_{\mathrm{runtime}} =
\begin{cases}
M_{\mathrm{override}} & \text{if an explicit override is given}, \\
M_{\mathrm{cuda}} + M_{\mathrm{alloc}} + M_{\mathrm{nccl}} + M_{\mathrm{ds}}
& \text{otherwise.}
\end{cases}
\]

\subsection{Persistent Backend Buffers}
One more resident term appears in some distributed full fine-tuning regimes:
persistent backend buffers. These are not model tensors in the usual sense.
They stand for backend-managed storage that remains live across phases. The
estimator keeps this term explicit because it is real memory, but it is not the
same kind of object as a parameter or an activation.

\subsection{The Two Resident Floors}
After collecting those pieces, the estimator builds two floors:
\begin{itemize}[leftmargin=1.5em]
\item the \textbf{base floor}, which includes everything resident except
gradients,
\item the \textbf{gradient floor}, which is the base floor plus gradients.
\end{itemize}

Written compactly:
\begin{align*}
R_{\mathrm{base}} &= M_{\mathrm{params}} + M_{\mathrm{opt}}
+ M_{\mathrm{master}} + M_{\mathrm{runtime}} + M_{\mathrm{persist}}, \\
R_{\mathrm{grad}} &= R_{\mathrm{base}} + M_{\mathrm{grads}}.
\end{align*}

These two floors are enough for the phase model:
\begin{itemize}[leftmargin=1.5em]
\item forward starts from $R_{\mathrm{base}}$,
\item backward and optimizer step start from $R_{\mathrm{grad}}$.
\end{itemize}

\section{Retained Activations}
This is the part of the model that is easiest to misread, because many people
informally use the word ``activations'' to mean several different things at
once.

\subsection{Visible Is Not The Same As Retained}
The estimator makes a deliberate distinction:
\begin{itemize}[leftmargin=1.5em]
\item \textbf{hook-visible outputs} are tensors you can see in instrumentation,
\item \textbf{retained forward state} is the subset that actually survives the
end of forward,
\item \textbf{backward activation state} is what remains alive during backward,
including backward-only context.
\end{itemize}

This distinction is not cosmetic. A hook can see a tensor that is later freed
before backward begins. Conversely, a tensor may be small in a hook trace but
have a large effect because it is replicated across layers or preserved through
checkpoint boundaries.

\subsection{The Retained-Forward Proxy}
The estimator gathers retained-forward state into one proxy term,
$M_{\mathrm{proxy}}$. It is built from named tensor classes rather than one
black-box multiplier. The important classes are:
\begin{itemize}[leftmargin=1.5em]
\item visible propagation state,
\item checkpoint-boundary inputs,
\item deduplicated saved linear inputs,
\item MLP intermediates when they must survive,
\item attention-saved state,
\item LoRA low-rank intermediates,
\item logits or loss context when needed,
\item widened attention-path contexts for architectures whose attention widths
are not all equal to the hidden size.
\end{itemize}

The key design choice here is explicit attribution. The proxy is not supposed
to answer ``how much forward stuff exists somewhere.'' It is supposed to answer
``which named forward objects are expected to survive into backward.''

The compact summary is:
\[
M_{\mathrm{proxy}} =
M_{\mathrm{visible}}
+ M_{\mathrm{saved\_inputs}}
+ M_{\mathrm{mlp}}
+ M_{\mathrm{attn}}
+ M_{\mathrm{lora}}
+ M_{\mathrm{checkpoint}}
+ M_{\mathrm{loss}}
+ M_{\mathrm{widened}}.
\]
Not every regime activates every term, but this is the right way to think about
the proxy: it is a sum of named retained object classes.

\paragraph{Full fine-tuning.}
When the whole stack is trainable and checkpointing is off, retained
activations grow roughly like:
\[
M_{\mathrm{proxy}}^{\mathrm{full\_ft,\ no\ ckpt}}
\propto
L \cdot T_{\mathrm{local}} \cdot
\left(H + I + d_{\mathrm{attn,path}}\right).
\]
That is the expensive regime most people expect: activation pressure grows
linearly with depth and with local token count, and the MLP as well as the
attention path can both matter.

\paragraph{LoRA.}
LoRA usually keeps much less activation state because only the targeted adapter
paths are trainable:
\[
M_{\mathrm{proxy}}^{\mathrm{lora,\ no\ ckpt}}
\propto
L \cdot T_{\mathrm{local}} \cdot
\left(d_{\mathrm{target\ path}} + r\right),
\]
where $r$ stands for the low-rank adapter width. The important point is not the
exact symbol choice. It is that LoRA growth is still linear in depth and token
count, but it is tied to the trainable adapter paths rather than the full
hidden-plus-MLP stack.

\paragraph{Backend contrast.}
If the only thing you want to remember is how eager and non-eager attention
differ, the comparison can be compressed into one small table:

\begin{center}
\small
\begin{tabular}{>{\raggedright\arraybackslash}p{1.0in}>{\raggedright\arraybackslash}p{1.8in}>{\raggedright\arraybackslash}p{2.65in}}
\toprule
Backend family & Extra retained-growth term & Interpretation \\
\midrule
eager (\texttt{standard}) &
$L \cdot T_{\mathrm{local}} \cdot S_{\mathrm{eff}} h_{\mathrm{local}}$ &
The estimator carries an attention-score-like state that grows with effective
sequence length and local head count. This is why eager attention becomes much
more sequence-sensitive. \\
non-eager (\texttt{sdpa}/\texttt{flash2}) &
$L_{\mathrm{act}} \cdot T_{\mathrm{local}} \cdot d_{\mathrm{attn,path}}$ &
The large score-retention term disappears. What remains is width-shaped context
from the attention path, so growth stays linear in tokens. \\
\bottomrule
\end{tabular}
\normalsize
\end{center}

Architecture still matters inside the non-eager family. The widened
attention-path terms make some models grow faster than a plain hidden-size
story would suggest:
\[
M_{\mathrm{widened}}
\propto
L_{\mathrm{act}} \cdot T_{\mathrm{local}} \cdot d_{\mathrm{widened\ path}}.
\]
So even when two models share the same hidden size, the one with wider
attention-path metadata can retain more activation state.

\subsection{What Checkpointing Changes}
Checkpointing does not magically remove all activation memory. What it changes
is the number of layers that are allowed to keep full internal state at once.

The current implementation uses a simple segmentation rule:
\begin{itemize}[leftmargin=1.5em]
\item without checkpointing, retained-forward classes can scale with the full
layer stack,
\item with checkpointing, only one checkpoint segment is allowed to retain its
full internal state at a time.
\end{itemize}

In the current code that effective segment length is one block. That keeps the
logic simple and aligns well with the measured checkpointed regimes in the
canonical corpus.

Viewed proportionally, checkpointing changes the layer factor much more than
the token factor:
\[
M_{\mathrm{proxy}}^{\mathrm{ckpt}}
\propto
L_{\mathrm{act}} \cdot T_{\mathrm{local}} \cdot d_{\mathrm{retained}},
\qquad
L_{\mathrm{act}} \ll L.
\]
So checkpointing does not make activation memory stop scaling with sequence
length. It mainly collapses the number of layers that can keep full retained
state at once.

Written compactly:
\[
L_{\mathrm{act}} =
\begin{cases}
L & \text{if checkpointing is off}, \\
1 & \text{if checkpointing is on}.
\end{cases}
\]

\subsection{Backward Activation State}
Backward activation memory is not rebuilt from scratch. The estimator starts
from the retained-forward proxy and then adds backward-only context:
\[
A_{\mathrm{bwd}} = M_{\mathrm{proxy}} + M_{\mathrm{bwd,extra}}.
\]

The extra backward term covers objects that do not need to survive the whole
forward pass but do need to exist while gradients are being computed. Examples
include residual or normalization context and some saved-input state that only
matters once backward begins.

Checkpointed distributed LoRA now uses a deliberately slim backward path. In
DDP, ZeRO-2, and ZeRO-3, the forward proxy is the checkpoint boundary plus the
loss-state term and the LoRA low-rank intermediates. Backward then adds only
the saved linear-input context and the residual or normalization context that
must become live once gradients are formed.

Viewed proportionally, backward keeps the same forward-retained base and adds a
smaller extra context term:
\[
A_{\mathrm{bwd}}
\propto
M_{\mathrm{proxy}}
+ T_{\mathrm{local}} \cdot d_{\mathrm{bwd,extra}}.
\]
This is why backward often remains larger than forward even after the retained
proxy has already been counted.

Earlier versions also added an extra ZeRO-3 checkpointed visible-propagation
slice on top of this path. Long-context OLMo 3 replay showed that this term
systematically overcounted backward activation state, so the current estimator
removes that addend and keeps the checkpointed distributed LoRA path aligned to
named retained objects only.

So the two activation views are:
\begin{align*}
A_{\mathrm{fwd}} &= M_{\mathrm{proxy}}, \\
A_{\mathrm{bwd}} &= M_{\mathrm{proxy}} + M_{\mathrm{bwd,extra}}.
\end{align*}

\subsection{Why The Estimator Reports Two Activation Views}
The public output exposes both:
\begin{itemize}[leftmargin=1.5em]
\item the retained-forward proxy, which is the best analytical match to the
measured retained-activation observable,
\item and the phase-local activation bucket that is used when a specific phase
wins the peak.
\end{itemize}

This separation keeps the calibration honest. A measured forward-retained proxy
should be compared to the estimator's retained-forward proxy, not to the
backward activation bucket and not to the public breakdown after a different
phase wins.
